\documentclass[10pt,aspectratio=169]{beamer}
\usetheme[progressstyle=movingCircCnt]{Feather}

\usepackage[utf8]{inputenc}
\usepackage[T1]{fontenc}
\usepackage[english,brazil]{babel}
\usepackage{helvet}
\usepackage{amsmath}
\usepackage{booktabs}
\usepackage{tikz}
\usetikzlibrary{arrows.meta,positioning}

\definecolor{unbblue}{RGB}{0,51,102}
\definecolor{unbgreen}{RGB}{33,140,33}
\definecolor{softgreen}{RGB}{230,244,230}
\setbeamercolor{alerted text}{fg=unbgreen}
\setbeamerfont{frametitle}{size=\large,series=\bfseries}
\setbeamertemplate{itemize item}{\color{unbgreen}$\blacktriangleright$}
\setbeamertemplate{itemize subitem}{\color{unbgreen}$\bullet$}

\title[Assinatura Digital RSA]{\textbf{Sistema de Assinatura Digital:\\RSA-OAEP e RSA-PSS}}
\subtitle[CIC0201]{Trabalho de Implementação 2 -- CIC0201}
\author[Grupo]{Ana Luísa Reis Nascente \and Gabriel de Sousa\\Marina Pimentel Moreno \and Pedro de Paula Campos}
\institute[]{Segurança Computacional\\Universidade de Brasília}
\date{29 de setembro de 2026}

\newcommand{\source}[1]{\vfill{\tiny\color{gray}Fonte: #1}}

\begin{document}

{\1
\begin{frame}[plain,noframenumbering]
  \titlepage
\end{frame}}

\section{Introdução}
\begin{frame}{Introdução}{RSA para cifragem $\times$ RSA para assinatura}
  % TODO
\end{frame}

\section{Parte I -- Chaves RSA}
\begin{frame}{Geração de chaves}{Miller-Rabin e parâmetros CRT}
  % TODO Gabriel
\end{frame}

\section{Parte II -- RSA-OAEP}
\begin{frame}{RSA-OAEP}{MGF1 e detecção de adulteração}
  % TODO Ana
\end{frame}

\section{Parte III -- RSA-PSS}
\begin{frame}{RSA-PSS}{Salt, MGF1 e Base64}
  % TODO Marina
\end{frame}

\section{Parte IV -- Verificação}
\begin{frame}{Verificação e adulteração}{Arquivo, assinatura e chave pública}
  % TODO Pedro
\end{frame}

\section{Parte V -- Análise de segurança}
\begin{frame}{Análise de segurança}{Textbook RSA, OAEP, PSS e Ed25519}
  % TODO Pedro
\end{frame}

\section{Demonstração}
\begin{frame}{Demonstração}
  \centering\large \texttt{bash exemplos/demo.sh}
\end{frame}

\section{Conclusão}
\begin{frame}{Conclusão}
  % TODO
\end{frame}

\end{document}
